The Large Hadron Collider (LHC) is currently the biggest collider in the world for hadrons. It facilitates measurements of pp and A+A collisions within 4 major experiments: ATLAS, ALICE, CMS and LHCb. The LHC had its major success after its first running at 2012 when the Higgs boson was discovered\cite{aad_observation_2012}. The collider is built out of multiple acceleration rings, which were used previously in LEP, etc. to accelerate protons and ions to TeV scale energies. A schematic of the accelerator complex can be seen in Fig. \ref{fig:LHC_compltex}.

\begin{figure}[h!] 
    \centering % Centers the image and caption
    \includegraphics[width=0.8\textwidth]{\main/illustrations/LHC_complex.png}
    \caption{A schema of the LHC complex with the injectors and different accelerator rings. Different facilities are color coded and dated \cite{lopienska_cern_2022}.}
    \label{fig:LHC_compltex}
\end{figure}

Heavy ion operations at the LHC are typically done once a year, and are usually designed around accelerating Pb ions, although other ions such as O, Ne and Xe have been run as special runs as well. 
